% ============================================
% II. RELATED WORK AND BACKGROUND
% ============================================
\section{Related Work and Background}
\label{sec:related}

Our research into optimizing blockchain-backed trading architectures draws upon several established domains of computer science and finance. We position our work at the intersection of blockchain scalability, concurrent system design, and high-frequency trading principles. The following sections review key literature in these areas and situate our contributions within the broader academic context.

\subsection{Blockchain Throughput and Scalability}

The performance limitations of public blockchains are well-documented. Initial quantitative analyses, such as the work by Bez et al.~\cite{bez2019scalability}, established a baseline for Ethereum's throughput, attributing its low transaction processing capacity to the inherent trade-offs between decentralization, security, and scalability. Their work provides a theoretical model for the bottlenecks our baseline architecture (Architecture~1) empirically demonstrates at 0.25 transactions per second (tx/sec). Our findings extend this analysis by showing that targeted architectural optimizations can achieve 31.19~tx/sec, a 2.8$\times$ improvement over conventional asynchronous designs, directly addressing the scalability challenge they identify. The full ablation from our deliberately naive baseline to the proposed architecture demonstrates a 125-fold cumulative improvement, though the more practically relevant comparison, against a conventional asynchronous multi-threaded design, yields a 2.8$\times$ gain.

More recent comparative studies, like the one conducted by Ucbas et al.~\cite{ucbas2023performance}, have benchmarked various blockchain platforms, including Ethereum and Hyperledger Fabric, under controlled loads. Their work highlights the significant performance differences between platforms. Our research complements these platform-level comparisons by providing a granular, application-level analysis. By systematically testing five distinct architectures for a trading system, we isolate specific bottlenecks, such as nonce contention and server-side locking, that are often abstracted away in broader platform evaluations. Our results, showing a latency range from 2.9s to 27.2s depending on the architectural design, underscore the critical role of application architecture in achieving high performance.

Layer~2 scaling solutions, such as rollups and sidechains, represent a popular and parallel avenue of research for enhancing blockchain throughput~\cite{song2024layer}. In blockchain terminology, Layer~1 refers to the base protocol and its consensus mechanism (e.g., Ethereum mainnet), while Layer~2 refers to secondary frameworks built atop the base layer to handle transactions off-chain before settling them on-chain. These Layer~2 approaches typically move computation and state storage off-chain to reduce the burden on the main network. Our work is orthogonal but complementary to this research thrust. We focus on optimizing the Layer~1 application architecture itself, which can be used in conjunction with Layer~2 solutions to achieve even greater scalability.

\subsection{Nonce Management and Transaction Contention}

A significant bottleneck we identified in our experiments is transaction contention arising from Ethereum's nonce mechanism. In Ethereum, a \textit{nonce} is a per-account transaction counter that ensures transactions are processed exactly once and in a deterministic order; each transaction must include a nonce exactly one greater than the previous confirmed transaction, and any gap or duplicate causes rejection. This mechanism, while essential for security and replay protection, requires that transactions from a single account be processed in a strict, sequential order. Our results for Architectures~2 and~3, where introducing multi-threading yielded only a 36\% throughput gain, empirically confirm that nonce contention is a primary limiting factor for parallelization. This finding aligns with theoretical models and security analyses of the Ethereum mempool, which highlight the challenges of high-volume transaction submission~\cite{wang2024mempool}.

To address this, some researchers have proposed alternative concurrency models. Paimani's work on SonicChain, for example, introduces a ``wait-free'' approach using concurrency delegation and pseudo-static annotations to mitigate conflicts~\cite{paimani2021sonicchain}. While this provides a path to reducing contention, our multi-wallet approach (Architectures~4 and~5) offers a more direct and aggressive partitioning strategy. By using multiple wallets, we effectively create independent nonce sequences, although, as our results for Architecture~4 show, this can shift the bottleneck to the server-side if not managed with appropriate concurrency controls.

\subsection{Concurrent Architectures for Financial Systems}

Our approach to resolving server-side contention draws heavily from principles developed in the high-frequency trading (HFT) domain. The catastrophic 27.2s latency observed in our unsynchronized multi-wallet architecture (Architecture~4) is a classic symptom of resource contention, such as cache thrashing and lock convoying, which HFT systems are meticulously designed to avoid. The foundational work on lock-free data structures by Moir and Shavit~\cite{moir2018concurrent} provides the theoretical underpinnings for the solution we implement in Architecture~5.

Specifically, we apply the principles of the LMAX Disruptor pattern, a high-performance inter-thread communication mechanism that uses a ring buffer and lock-free operations to achieve extremely low latency~\cite{thompson2011disruptor}. By implementing a symbol-sharded, lock-free queueing system, we were able to reduce the average transaction latency from 27.2s down to 15.4s, a 44\% improvement. This empirically validates that techniques proven in traditional finance, as detailed in works by Bilokon and Gunduz~\cite{bilokon2023cpp} and Ng and Malik~\cite{ng2018lockfree}, can be successfully adapted to blockchain-based systems to overcome critical performance barriers.

\subsection{Blockchain-Based Trading and Decentralized Finance}

The architecture of decentralized exchanges (DEXs) has evolved significantly, from on-chain order books to automated market makers (AMMs)~\cite{xu2023sok}. However, many existing designs suffer from scalability limitations and are susceptible to fairness issues like Maximal Extractable Value (MEV), where miners or validators reorder transactions for their own financial gain~\cite{alipanahloo2024mev}.

Our symbol-sharded architecture directly addresses these challenges. By creating isolated, per-symbol processing queues, we not only improve scalability but also enhance fairness. Within each shard, transactions are processed in a deterministic, first-in-first-out (FIFO) manner, which inherently mitigates the potential for MEV strategies like front-running and sandwich attacks that rely on arbitrary transaction reordering. Recent analyses of the Ethereum blockchain have shown that such favoritism in transaction ordering is a systemic issue~\cite{mancino2025mev}. Our work provides a concrete architectural solution that promotes a more equitable and predictable trading environment.

\subsection{Sharding and Data Partitioning}

Sharding is a well-established technique for scaling distributed databases and is now being actively developed for blockchains like Ethereum~2.0. Early research in this area, such as the work by Fynn and Pedone~\cite{fynn2018partitioning}, explored various methods for partitioning blockchain state and highlighted the significant overhead associated with cross-shard communication. Our approach leverages a domain-specific partitioning strategy, sharding by trading symbol, which is more efficient for a trading workload than generic graph-based partitioning methods.

Our work is also highly complementary to research on optimizing cross-shard communication. Kudzin et al.~\cite{kudzin2022sharding} have proposed novel data structures to compress transaction data in rollups, thereby increasing the number of cross-shard transactions that can fit into a single block. While their work focuses on the communication between shards, our work optimizes the processing \textit{within} a shard. A system that combines both our symbol-sharded architecture and their cross-shard compression techniques could achieve a new level of performance and scalability for decentralized trading.

\subsection{Research Gap and Motivation}

The literature reviewed above reveals several critical gaps that directly motivate our research questions. First, although lock-free and HFT-inspired architectures are well-established in traditional finance, their adaptation to the unique constraints of blockchain transaction submission (particularly nonce sequencing) remains under-explored, motivating \textbf{RQ2}. Second, while blockchain scalability has been extensively studied at the platform level, there is a lack of systematic, application-level analysis that isolates specific bottlenecks such as nonce contention and server-side locking, the focus of \textbf{RQ1}. Finally, while latency is recognized as critical in trading systems, the relationship between architectural design choices in the submission layer and end-to-end settlement latency has not been empirically characterized for blockchain systems, which is the focus of \textbf{RQ3}. This paper addresses these gaps through a systematic ablation study that isolates each factor.

